\chapter{概率：基础与大数
定律}\label{ch:b3:probability}

本科二年级在可数空间上建立了概率；\omterm{def:b3:measure:measure}{测度}论
现在去掉了一切限制。\omterm{def:b3:probability:space}{概率空间}是
总\omterm{def:b3:measure:measure}{测度}为 $1$ 的\omterm{def:b3:measure:measure}{测度空间}，\omterm{def:b3:probability:space}{随机变量}是
可测映射，\omterm{def:b3:probability:space}{期望}是勒贝格积分——并且整套分析武器库
一并
（\cref{ch:b3:measure,ch:b3:lebesgue,ch:b3:product}）适用于随机现象。本章建立
这套对应关系，构造
\omterm{def:b3:probability:independence}{独立}\omterm{def:b3:probability:space}{随机变量}的无穷序列（在
$\intcc01$ 上，由二进制数字给出：随机性就藏在
\omterm{def:b3:measure:lebesgueouter}{勒贝格测度}之中），证明博雷尔--坎泰利引理，
柯尔莫哥洛夫零一律，厘清收敛的各种
方式，并证明大数定律——这一定理使
频率收敛于概率，使统计成为可能。
周末问题给出埃特马迪
关于强大数定律在最终 $L^1$ 形式下的证明。

\section{对应关系}

\begin{definition}\label{def:b3:probability:space}
\emph{概率空间}\index{概率空间}是一个
\omterm{def:b3:measure:measure}{测度空间} $(\Omega, \mathcal A, \P)$，满足 $\P(\Omega) =
1$；$\mathcal A$ 中的元素称为\emph{事件}，一个性质
\emph{几乎必然}（a.s.）成立，是指其对应事件的
概率为 $1$。\emph{随机变量}\index{随机
变量}是可测映射 $X \colon \Omega \to \R$（或
$\R^d$：随机向量）；其\emph{分布}\index{分布}
是推前概率\omterm{def:b3:measure:measure}{测度} $\P_X =
X_*\P$ 在 $\R$ 上（\cref{exo:b3:product:9}），由
\emph{分布函数} $F_X(t) = \P(X \leq t)$ 确定
（\cref{exo:b3:measure:3}）。$X$ 具有\emph{密度} $f$，若
$\P_X = f\,\dd\lambda$；它为\emph{离散}的，若 $\P_X$ 是
可数个狄拉克\omterm{def:b3:measure:measure}{测度}的组合。
\emph{期望}\index{期望}为
\[
\E[X] = \int_\Omega X\,\dd\P
\qquad (X \geq 0 \text{ 或 } X \in L^1(\P)),
\]
而\emph{转换定理}（\cref{exo:b3:product:9}）
在分布意义下计算：$\E[g(X)] = \int_\R g\,\dd\P_X$——
$= \sum g(x_k)p_k$ 为离散情形，$= \int
g(x)f(x)\dd x$ 为密度情形：本科二年级的公式，如今
已成为同一理论的定理。\emph{方差}为 $\V(X) =
\E[(X - \E X)^2] = \E[X^2] - (\E X)^2$ 对 $X \in L^2$ 成立。
\end{definition}

\begin{example}\label{ex:b3:probability:laws}
标准分布及其值得注意的性质：伯努利
分布 $\mathcal B(p)$、二项分布 $\mathcal B(n, p)$、几何分布，
泊松分布 $\mathcal P(\lambda)$（离散情形，本科二年级的表
格仍然成立）；$\intcc01$ 上的均匀分布（即\omterm{def:b3:measure:lebesgueouter}{勒贝格测度}
本身）；指数分布 $\mathcal E(\lambda)$（密度
$\lambda\eu^{-\lambda x}\mathbf 1_{x>0}$）；
以及\emph{高斯分布}\index{高斯分布} $\mathcal N(m, \sigma^2)$
密度为 $\frac1{\sigma\sqrt{2\pi}}\exp\bigl(-\frac{(x -
m)^2}{2\sigma^2}\bigr)$——是概率密度，由
\cref{pb:b3:lebesgue:1}可知，其均值为 $m$，方差为
$\sigma^2$（高斯矩，见\cref{exo:b3:product:10}）。
\end{example}

\begin{proposition}[马尔可夫与切比雪夫]
\label{prop:b3:probability:markov}
对 $X \geq 0$ 及 $a > 0$，有 $\P(X \geq a) \leq \frac{\E
X}{a}$；对 $X \in L^2$，有 $\P\bigl(\abs{X - \E X} \geq
a\bigr) \leq \frac{\V(X)}{a^2}$。\index{马尔可夫
不等式}\index{切比雪夫不等式}
\end{proposition}

\begin{proof}
由\cref{exo:b3:lebesgue:5}(a)可知马尔可夫不等式成立；切比雪夫不等式即将马尔可夫
不等式用于 $(X - \E X)^2$。
\end{proof}

\section{独立性}

\begin{definition}\label{def:b3:probability:independence}
子 $\sigma$-代数 $\mathcal A_1, \dots, \mathcal A_n
\subseteq \mathcal A$ 称为\emph{独立}\index{独立}的，
若 $\P(A_1\cap\dots\cap A_n) = \prod\P(A_i)$ 对所有 $A_i \in
\mathcal A_i$；事件独立，是指其 $\sigma$-代数
$\{\varnothing, A_i, A_i^c, \Omega\}$ 独立；\omterm{def:b3:probability:space}{随机变量}
$X_1, \dots, X_n$ 独立，若 $\sigma$-代数 $\sigma(X_i) =
X_i^{-1}(\mathcal B(\R))$ 皆独立。称无穷族
为独立，若每个有限子族都独立。
\end{definition}

\begin{theorem}\label{thm:b3:probability:independence}
$X_1, \dots, X_n$ \omterm{def:b3:probability:independence}{独立}，当且仅当向量
$(X_1, \dots, X_n)$ 的分布为\omterm{thm:b3:product:existence}{乘积测度}
$\P_{X_1}\otimes\cdots\otimes\P_{X_n}$。此时，对
$g_i \geq 0$（或使乘积\omterm{def:b3:lebesgue:l1}{可积}的）：
\[
\E\Bigl[\prod_ig_i(X_i)\Bigr] = \prod_i\E[g_i(X_i)],
\]
特别地，$\E[XY] = \E X\,\E Y$ 及 $\V(X_1 + \dots +
X_n) = \sum\V(X_i)$ 对\omterm{def:b3:probability:independence}{独立}的 $L^2$ 变量成立。
\end{theorem}

\begin{proof}
若 $X_i$ \omterm{def:b3:probability:independence}{独立}，则两个概率\omterm{def:b3:measure:measure}{测度}
$\P_{(X_1,\dots,X_n)}$ 与 $\bigotimes\P_{X_i}$ 在所有
博雷尔集的乘积 $B_1\times\dots\times B_n$ 上一致——一个
$\pi$-系，它生成 $\mathcal B(\R^n)$
（见\cref{prop:b3:product:sections}(b)）——从而处处
（见\cref{thm:b3:measure:uniqueness}）。反之，乘积分布
将一切事件 $\bigcap_iX_i^{-1}(B_i)$ 分解，即得\omterm{def:b3:probability:independence}{独立}性。\omterm{def:b3:probability:space}{期望}公式于是
由托内利/富比尼（\cref{thm:b3:product:tonelli}）经
转换定理得到；$\E[XY] = \E X\E Y$ 即取 $g_i =
\mathrm{id}$ 的情形，展开平方即得方差的
可加性（交叉项 $\E[(X_i - \E X_i)(X_j - \E X_j)]
= 0$）。
\end{proof}

\begin{theorem}[独立序列的存在性]
\label{thm:b3:probability:existence}
在 $\bigl(\intcc01, \mathcal L, \lambda\bigr)$ 上存在
\omterm{def:b3:probability:independence}{独立}\omterm{def:b3:probability:space}{随机变量}序列 $(U_n)_{n\geq1}$，
每个都在 $\intcc01$ 上均匀分布。由此，对任意给定的
分布 $(\mu_n)$（定义在 $\R$ 上），存在相互\omterm{def:b3:probability:independence}{独立}的 $(X_n)$，使得
$\P_{X_n} = \mu_n$。\index{独立序列的存在性}
\end{theorem}

\begin{proof}
\emph{二进制数字。} 对 $\omega \in \intcc01$，记 $(b_k(\omega))$
为其二进制数字（$\omega = \sum b_k2^{-k}$；取
不以全为 $1$ 的尾部结束的展开；歧义仅涉及
一个可数、因而零测的集合）。每个 $b_k$ 是\omterm{def:b3:probability:space}{随机
变量}（$\{b_k = 1\}$ 是有限个二进
区间之并），且向量 $(b_1, \dots, b_m)$ 取遍
$\{0,1\}^m$ 中的每个值都在长为 $2^{-m}$ 的二进区间上取得：
各 $b_k$ 相互\omterm{def:b3:probability:independence}{独立}，均服从伯努利$(\frac12)$分布。

\emph{重新分组。} 把 $\N^*$ 分成无穷多个两两不交的
无限集 $(I_n)$（例如按素数幂，或按对角线）；
令 $(k^n_j)_j$ 枚举 $I_n$，并记
\[
U_n = \sum_{j\geq1} b_{k^n_j}\,2^{-j} .
\]
每个 $U_n$ 服从均匀分布：其二进制数字是\omterm{def:b3:probability:independence}{独立}的
公平比特，故 $\P(U_n \in [l2^{-m}, (l+1)2^{-m})) = 2^{-m}$ 对
每个二进区间成立，且二进区间决定分布
（见\cref{thm:b3:measure:uniqueness}）。$U_n$ 相互\omterm{def:b3:probability:independence}{独立}：
它们是\omterm{def:b3:probability:independence}{独立}族 $(b_k)$ 中互不相交的块的函数——
形式地说，事件 $\{U_n \in D_n\}$（$D_n$ 为
二进区间）仅依赖于来自互不相交集合的有限多个数字，
从而可分解；再由 $\pi$-系论证推广到
一切博雷尔集。

\emph{任意分布。} 设 $G_n(u) = \inf\{t : F_{\mu_n}(t)
\geq u\}$（即分布函数 $F_{\mu_n}$ 的\emph{分位函数}）；
关键等价关系 $G_n(u) \leq t
\iff u \leq F_{\mu_n}(t)$（由 $F$ 的右连续性，
及单调性）可知 $X_n = G_n(U_n)$ 可测，且
$\P(X_n \leq t) = \P(U_n \leq F_{\mu_n}(t)) =
F_{\mu_n}(t)$，即分布为 $\mu_n$；\omterm{def:b3:probability:independence}{独立}性随之继承
（\omterm{def:b3:probability:independence}{独立}变量的函数，
见\cref{exo:b3:probability:3}）。
\end{proof}

\begin{example}[生日问题，如实计算]
\label{ex:b3:probability:birthday}
在 $n$ 个人的生日相互\omterm{def:b3:probability:independence}{独立}、且在
$N = 365$ 天上均匀分布的条件下，所有生日互不相同的概率为
\[
p_n = \prod_{k=1}^{n-1}\Bigl(1 - \frac kN\Bigr),
\]
可由逐次条件化得到（或直接：有利情形数
$N(N-1)\cdots(N - n + 1)$ 比上总数 $N^n$，这一计数论证
借助\omterm{def:b3:probability:independence}{独立}性的乘积公式变得严格）。
取对数，并利用 $-\ln(1 - x) = x +
O(x^2)$：
\[
\ln p_n = -\frac{n(n-1)}{2N} +
O\Bigl(\frac{n^3}{N^2}\Bigr),
\qquad\text{故}\qquad
p_n \approx \eu^{-n^2/2N} .
\]
临界点 $p_n = \frac12$ 位于 $n \approx
\sqrt{2N\ln2} \approx 1.18\sqrt N$ 处：取 $N = 365$ 时 $n = 23$
（$p_{23} = 0.4927$）。两点启示。其一，$n$ 个
元素落入 $N$ 个盒子时的碰撞出现在 $n \sim \sqrt N$ 的量级上，
而非 $n \sim N$——这正是支配哈希碰撞的\emph{生日尺度}，
以及密码学中生日攻击的 $\sqrt N$ 代价。其二，
这一计算是一个模板：
各 $\binom n2$ 个两两碰撞事件并不\omterm{def:b3:probability:independence}{独立}，
然而答案的表现却仿佛它们是\omterm{def:b3:probability:independence}{独立}的——
（$\eu^{-\binom n2/N}$ 恰好就是\omterm{def:b3:probability:independence}{独立}配对的
启发式近似）——泊松近似的一个初步实例，
其严格化见\cref{ch:b3:clt}中的周末问题
（勒卡姆不等式）。
\end{example}

\section{博雷尔--坎泰利与零一律}

\begin{theorem}[博雷尔--坎泰利]
\label{thm:b3:probability:borelcantelli}
设 $(A_n)$ 为事件列，并记 $\limsup A_n = \bigcap_N
\bigcup_{n\geq N}A_n$（“$A_n$ 无穷多次
发生”）。\index{博雷尔--坎泰利引理}
\begin{enumerate}
  \item 若 $\sum\P(A_n) < \infty$，则 $\P(\limsup A_n) =
        0$。
  \item 若 $\sum\P(A_n) = \infty$ \emph{且各 $A_n$
        相互\omterm{def:b3:probability:independence}{独立}}，则 $\P(\limsup A_n) = 1$。
\end{enumerate}
\end{theorem}

\begin{proof}
(1) 由\cref{exo:b3:measure:4}可知。(2)：对 $N \leq M$，
由余事件的\omterm{def:b3:probability:independence}{独立}性（见\cref{exo:b3:probability:3}）
可得
\[
\P\Bigl(\bigcap_{n=N}^{M}A_n^c\Bigr) = \prod_{n=N}^M\bigl(1
- \P(A_n)\bigr) \leq
\exp\Bigl(-\sum_{n=N}^M\P(A_n)\Bigr) \xrightarrow[M \to
\infty]{} 0
\]
（$1 - x \leq \eu^{-x}$；级数发散）。
因此 $\P\bigl(\bigcup_{n\geq N}A_n\bigr) = 1$ 对每个 $N$ 成立，而
关于 $N$ 的递减交集的概率仍为
$1$（由上连续性，见\cref{prop:b3:measure:basics}）。
\end{proof}

\begin{theorem}[柯尔莫哥洛夫 0--1 律]
\label{thm:b3:probability:zeroone}
设 $(X_n)$ \omterm{def:b3:probability:independence}{独立}，记 $\mathcal T =
\bigcap_N\sigma(X_N, X_{N+1}, \dots)$ 称为\emph{尾
$\sigma$-代数}\index{尾 $\sigma$-代数}（对有限个
$X_n$ 的任意改变都不敏感的事件：$\sum X_n$ 的收敛性、
$\frac{S_n}n$ 的收敛性、$\limsup$ 的取值，\dots）。
则每个 $T \in \mathcal T$ 都满足 $\P(T) \in \{0,
1\}$。\index{0--1 律}
\end{theorem}

\begin{proof}
固定 $N$。$\sigma$-代数 $\sigma(X_1, \dots, X_N)$ 与
$\sigma(X_{N+1}, \dots)$ \omterm{def:b3:probability:independence}{独立}：依赖于不相交变量块的事件，
在相应的生成 $\pi$-系上乘积分解
（柱集 $\bigcap_{i\leq N}\{X_i \in B_i\}$，分别为
关于后续变量的有限条件），再由邓金
（\cref{thm:b3:measure:dynkin}，两次应用，
每次一侧）将分解式推广。尾事件 $T$ 属于
$\sigma(X_{N+1}, \dots)$（对每个 $N$）：$T$ 与
每个 $\sigma(X_1, \dots, X_N)$ \omterm{def:b3:probability:independence}{独立}，故与
它们生成的 $\sigma$-代数 $\sigma(X_1, X_2, \dots)$ \omterm{def:b3:probability:independence}{独立}（再用一次邓金：
各 $\sigma(X_1,\dots,X_N)$ 的并是生成它的 $\pi$-系）。
但 $T \in \sigma(X_1, X_2, \dots)$ 也成立：$T$
与\emph{自身}\omterm{def:b3:probability:independence}{独立}，$\P(T) = \P(T\cap T) =
\P(T)^2$，故 $\P(T) \in \{0, 1\}$。
\end{proof}

\section{收敛方式}

\begin{definition}\label{def:b3:probability:modes}
$X_n \to X$ \emph{几乎必然收敛}，若 $\P(X_n \to X) = 1$；
\emph{依概率收敛}，若 $\P(\abs{X_n - X} \geq \varepsilon)
\to 0$ 对每个 $\varepsilon > 0$ 成立；\emph{$L^p$ 收敛}，若
$\E\abs{X_n - X}^p \to 0$。\index{随机变量的
收敛}
\end{definition}

\begin{proposition}\label{prop:b3:probability:modes}
(a) 几乎必然收敛蕴含依概率收敛；
(b) $L^p$ 收敛蕴含依概率收敛；
(c) 依概率收敛蕴含
\emph{沿某一子列}几乎必然收敛；
(d) 一般而言不存在其他蕴含关系。
\end{proposition}

\begin{proof}
(a) $\P(\abs{X_n - X} \geq \varepsilon) \leq
\P\bigl(\sup_{m\geq n}\abs{X_m - X} \geq \varepsilon\bigr)
\downarrow \P\bigl(\limsup\{\abs{X_m - X} \geq
\varepsilon\}\bigr) = 0$（几乎必然收敛下成立，
由上连续性；limsup 事件排除了收敛）。(b)
马尔可夫不等式：$\P(\abs{X_n - X} \geq \varepsilon) \leq
\varepsilon^{-p}\,\E\abs{X_n - X}^p$。(c) 取 $n_k$，使
$\P(\abs{X_{n_k} - X} \geq 2^{-k}) \leq 2^{-k}$；
由博雷尔–坎泰利 (1)，$\abs{X_{n_k} - X} < 2^{-k}$
最终几乎必然成立。(d) 打字机序列（见\cref{exo:b3:lp:3}）在
$(\intcc01, \lambda)$ 上依 $L^1$ 收敛且依概率收敛，
但处处不逐点收敛；$n\mathbf 1_{\intoo0{1/n}} \to 0$
几乎必然收敛，但不在 $L^1$ 中收敛；细节及其余反例见
\cref{exo:b3:probability:6}。
\end{proof}

\section{大数定律}

全文 $(X_n)$ \omterm{def:b3:probability:independence}{独立}同分布
（\emph{i.i.d.}），记 $S_n = X_1 + \dots + X_n$。

\begin{theorem}[弱大数定律]
\label{thm:b3:probability:wlln}
若 $X_1 \in L^2$，记 $m = \E X_1$：
\[
\P\Bigl(\Bigl|\frac{S_n}{n} - m\Bigr| \geq
\varepsilon\Bigr) \leq
\frac{\V(X_1)}{n\,\varepsilon^2}
\xrightarrow[n\to\infty]{} 0 :
\]
$\frac{S_n}n \to m$ 依概率收敛（且在 $L^2$ 中收敛）。\index{弱大数
定律}
\end{theorem}

\begin{proof}
$\E\frac{S_n}n = m$，且 $\V\bigl(\frac{S_n}n\bigr) =
\frac{n\V(X_1)}{n^2}$
由\cref{thm:b3:probability:independence}及切比雪夫不等式可得。
\end{proof}

\begin{theorem}[强大数定律]
\label{thm:b3:probability:slln}
若 $X_1 \in L^1$，则\index{强大数定律}
\[
\frac{S_n}{n} \xrightarrow[n\to\infty]{\text{几乎必然}} \E[X_1].
\]
此处仅在更强的假设 $X_1 \in
L^4$ 下证明；一般情形（$L^1$，埃特马迪的证明）为
周末问题。
\end{theorem}

\begin{proof}[在 $\E X_1^4 < \infty$ 下的证明]
中心化（$X_i \mapsto X_i - m$），不妨设 $m = 0$。展开：
\[
\E[S_n^4] = \sum_{i,j,k,l}\E[X_iX_jX_kX_l]
= n\,\E[X_1^4] + 3n(n-1)\,\bigl(\E[X_1^2]\bigr)^2 \leq
C\,n^2 ,
\]
由于\omterm{def:b3:probability:independence}{独立}性与中心化，所有含
孤立因子的项皆为零（$\E[X_iX_jX_kX_l] =
\E[X_i]\E[\cdots] = 0$，除非指标两两配对：唯一
幸存的只有 $n$ 项 $i=j=k=l$，以及
$3n(n-1)$ 个含两对不同指标的项）。马尔可夫不等式：
\[
\P\Bigl(\Bigl|\frac{S_n}n\Bigr| \geq \varepsilon\Bigr)
= \P\bigl(S_n^4 \geq n^4\varepsilon^4\bigr)
\leq \frac{Cn^2}{n^4\varepsilon^4} =
\frac{C}{\varepsilon^4n^2},
\]
该级数收敛。由博雷尔–坎泰利 (1)，对每个有理数
$\varepsilon$，$\abs{S_n/n} < \varepsilon$ 最终成立，
几乎必然；对 $\varepsilon \in \Q_+^*$ 取交（
可数个概率为 $1$ 的事件）：$S_n/n \to 0$ 几乎必然。
\end{proof}

\begin{example}[强大数定律能带来什么]
\label{ex:b3:probability:sllnapps}
(a) \emph{频率}：对\omterm{def:b3:probability:independence}{独立}同分布的掷硬币试验，观测到的
正面出现的频率几乎必然收敛于 $p$——这正是
概率本身的经验依据。
(b) \emph{蒙特卡罗}\index{蒙特卡罗}：设 $g \in
L^1(\intcc01)$，$(U_n)$ 为\omterm{def:b3:probability:independence}{独立}同分布的均匀\omterm{def:b3:probability:space}{随机变量}
（见\cref{thm:b3:probability:existence}），
$\frac1n\sum_{k\leq n}g(U_k) \to \int_0^1g$ 几乎必然成立：通过抽样计算积分，
可用于任意维数，与维数无关的速率 $\sim n^{-1/2}$
在\cref{ch:b3:clt}中得到精确表述。
(c) \emph{正规数}：几乎所有实数的二进制展开中，
“1”出现的渐近频率为 $\frac12$，
（将强大数定律用于
\cref{thm:b3:probability:existence}中的数字变量）——博雷尔定理，一个
关于\emph{每个}数的陈述，却由\omterm{def:b3:measure:measure}{测度}证明：
\cref{pb:b3:probability:1}将其推广到所有进制。
\end{example}

\begin{method}\label{met:b3:probability:toolkit}
关于随机序列渐近性陈述的工作步骤如下：
(1) \emph{该事件是否为尾事件？}若是，
则其概率为 $0$ 或 $1$，
（见\cref{thm:b3:probability:zeroone}），只需
判断是哪一个。(2) \emph{证明几乎必然成立的陈述}：
博雷尔–坎泰利——对“坏”事件，概率要求可和，
借助马尔可夫型、切比雪夫型界估计，
对已存在的各阶矩进行估计；\omterm{def:b3:probability:independence}{独立}性仅在
逆方向中才需要。(3) \emph{子列 + 夹逼}：先证明
沿某个易处理的子列收敛；再控制
其间的振荡，借助单调性或极大
不等式——这正是埃特马迪证明的骨架。(4) 对
分布极限，则在\cref{ch:b3:clt}中再讨论。
\end{method}

\section{练习}

\begin{exercise}[$\star$]\label{exo:b3:probability:1}
(a) 设 $X$ 的分布函数 $F$ 连续且严格递增，
证明 $F(X)$ 在 $\intcc01$ 上服从均匀分布，
且 $G(U) \sim F$，这里 $U$ 服从均匀分布，$G = F^{-1}$：
即通过求逆进行模拟。
(b) 求 $X^2$ 的分布函数与密度，
其中 $X$ 在 $\intcc{-1}1$ 上服从均匀分布，以及 $-\frac1\lambda\ln
U$ 的分布函数与密度，其中 $U$ 在 $\intoo01$ 上服从均匀分布。
\end{exercise}

\begin{exercise}[$\star$]\label{exo:b3:probability:2}
(a) 借助变换定理，计算泊松分布 $\mathcal
P(\lambda)$ 与几何分布的均值与方差。
(b) 证明：若正\omterm{def:b3:probability:space}{随机变量} $T$ 满足 $\P(T > t)
> 0$ 对所有 $t$ 成立，并具有\emph{无记忆}性质 $\P(T > t + s \mid
T > t) = \P(T > s)$ 对所有 $s, t \geq 0$ 成立，当且仅当 $T$ 服从
指数分布。\emph{（生存函数满足柯西
函数方程；单调性代替连续性。）}
\end{exercise}

\begin{exercise}[$\star\star$]\label{exo:b3:probability:3}
(a) 证明：若 $X_1, \dots, X_n$ 相互\omterm{def:b3:probability:independence}{独立}，且 $f_i$
为博雷尔函数，则 $f_i(X_i)$ 相互\omterm{def:b3:probability:independence}{独立}。
(b) 证明：事件 $A_1, \dots, A_n$ 相互\omterm{def:b3:probability:independence}{独立}，当且仅当
它们的余事件相互\omterm{def:b3:probability:independence}{独立}，当且仅当示性函数 $\mathbf 1_{A_i}$
为相互\omterm{def:b3:probability:independence}{独立}的\omterm{def:b3:probability:space}{随机变量}。
(c) （两两\omterm{def:b3:probability:independence}{独立}是较弱的条件）两枚均匀硬币：$A =$ 第一枚
正面朝上，$B =$ 第二枚正面朝上，$C =$ 两枚结果相同。证明
$A, B, C$ 两两\omterm{def:b3:probability:independence}{独立}，但不相互\omterm{def:b3:probability:independence}{独立}。
\end{exercise}

\begin{exercise}[$\star\star$]\label{exo:b3:probability:4}
(a) （无限猴子）\omterm{def:b3:probability:independence}{独立}同分布的均匀击键序列，
在有限字母表上几乎必然包含每个有限文本
无穷多次：用博雷尔–坎泰利 (2)，对
互不相交的块加以证明。
(b) （游程）对\omterm{def:b3:probability:independence}{独立}同分布的公平比特，记 $R_n$ 为
从位置 $n$ 开始的连续 1 段长度。证明几乎必然
$R_n \geq (1+\varepsilon)\log_2n$ 仅出现有限次，且 $R_n
\geq \log_2 n$ 出现无穷多次 \emph{（博雷尔–坎泰利的两半；
对后者，取互不相交的块以
获得\omterm{def:b3:probability:independence}{独立}性）}：前 $n$ 位数字中最长的游程
增长阶为 $\log_2n$。
\end{exercise}

\begin{exercise}[$\star\star$]\label{exo:b3:probability:5}
设 $(X_n)$ 相互\omterm{def:b3:probability:independence}{独立}。
(a) 证明 $\sum X_n z^n$ 的收敛半径
几乎必然为常数（可能为 $0$ 或 $\infty$）。
(b) 证明 $\P(\sum X_n \text{ 收敛}) \in \{0, 1\}$
以及 $\P(S_n/n \to m) \in \{0,1\}$。
(c) 给出一个关于 $(X_n)$ 的事件，它\emph{不是}尾事件，
并验证0–1律对它可能不成立。
\end{exercise}

\begin{exercise}[$\star\star$]\label{exo:b3:probability:6}
在 $(\intcc01, \lambda)$ 上——并附证明——给出\omterm{def:b3:probability:space}{随机变量}，使其
满足：
(a) $X_n \to 0$ 依概率收敛，且在每个 $L^p$ 中收敛，但
几乎必然处处不收敛；
(b) $X_n \to 0$ 几乎必然成立，但在任何 $L^p$ 中都不收敛；
(c) $X_n \to 0$ 在 $L^1$ 中收敛，但在 $L^2$ 中不收敛；
(d) 并证明：若 $X_n \to X$ 依概率成立，且 $\abs{X_n}
\leq Y \in L^1$，则 $X_n \to X$ 在 $L^1$ 中收敛
\emph{（子列 + 控制收敛 +
子列的子列技巧）}。
\end{exercise}

\begin{exercise}[$\star\star$]\label{exo:b3:probability:7}
一项民意调查要估计未知比例 $p$，
以经验频率 $\hat p_n$ 估计，其中 $n$ 为\omterm{def:b3:probability:independence}{独立}抽样次数。
(a) 用切比雪夫不等式证明 $\P(\abs{\hat p_n - p} \geq \varepsilon)
\leq \frac1{4n\varepsilon^2}$ （利用 $p(1-p) \leq \frac14$）。
(b) 按此界，至少需抽取多少次，才能保证误差 $\leq 3\%$ 的
概率 $\geq 95\%$？（真正答案
借助\cref{ch:b3:clt}，约为 $1070$，切比雪夫
界诚实却粗糙。）
\end{exercise}

\begin{exercise}[$\star\star\star$]\label{exo:b3:probability:8}
（伯恩斯坦）设 $f \in \mathcal C(\intcc01)$，定义
伯恩斯坦多项式 $B_nf(x) =
\sum_{k=0}^n\binom nkx^k(1-x)^{n-k}f\bigl(\frac kn\bigr)$.
(a) 识别 $B_nf(x) = \E\bigl[f\bigl(\frac
{S_n}n\bigr)\bigr]$，其中 $S_n$ 服从二项分布 $\mathcal B(n, x)$。
(b) 证明 $B_nf \to f$ 在 $\intcc01$ 上\emph{一致收敛}：
对事件 $\{\abs{\frac{S_n}n - x} \leq \delta\}$ 及其
补集分别讨论，利用一致连续性以及切比雪夫不等式，
并使用一致界 $\V(\frac{S_n}n) \leq \frac1{4n}$。
(c) 由此得出结论：这是魏尔斯特拉斯逼近定理的第二种
概率论式证明
（见\cref{cor:b3:complete:weierstrass}），并给出显式的收敛速度
$\norm{B_nf - f}_\infty \leq \frac32\,\omega_f(n^{-1/2})$
其中 $\omega_f$ 为连续模——至少证明
$O(\omega_f(n^{-1/2}))$ 形式的结论。
\end{exercise}

\begin{exercise}[$\star\star\star$]\label{exo:b3:probability:9}
（赠券收集者问题）从 $n$ 类卡片中均匀地抽取，
有放回地进行。设 $T_n$ 为直到所有
类型都出现过为止的抽取次数。
(a) 写出 $T_n = \sum_{k=1}^{n}\tau_k$，其中 $\tau_k$
服从参数为 $\frac{n - k + 1}n$ 的几何分布，各 $\tau_k$
相互\omterm{def:b3:probability:independence}{独立}，由此推出 $\E T_n = n\,H_n \sim n\ln n$
（$H_n$ 为调和数）及 $\V(T_n) \leq
\frac{\pi^2}6n^2$。
(b) 由切比雪夫不等式，$\frac{T_n}{n\ln n} \to 1$ 依概率成立。
(c) 用博雷尔–坎泰利引理细化：直接证明 $\P(T_n >
\beta n\ln n) \leq n^{1 - \beta}$，其中 $\beta > 1$，
\emph{（对“某类型在
$\beta n\ln n$ 次抽取后仍未被抽到”这一事件使用并集界，并利用 $1 - x \leq \eu^{-x}$）}，
由此推出：沿 $n = 2^m$，几乎必然 $T_n \leq \beta n\ln
n$ 最终成立，对每个 $\beta > 2$。
\end{exercise}

\begin{exercise}[$\star\star$]\label{exo:b3:probability:10}
利用数字构造
（见\cref{thm:b3:probability:existence}）：
(a) 直接计算验证 $U = \sum b_{2k}2^{-k}$
（取均匀随机数 $\omega$ 的偶数位数字）服从均匀分布，
且与 $V = \sum b_{2k-1}2^{-k}$ \omterm{def:b3:probability:independence}{独立}；
(b) 推出 $\intcc01$ 与
$\intcc01^2$ 之间存在一个去掉零测集后为双射的保测可测映射，并作评注：
一个均匀随机数中包含两个（乃至可数多个）
相互\omterm{def:b3:probability:independence}{独立}的随机数——可与皮亚诺曲线作比较，
（见\cref{pb:b3:topology:1}），后者实现了满射，但
不保测，也不是单射。
\end{exercise}

\begin{exercise}[$\star\star$]\label{exo:b3:probability:11}
（记录）设 $(X_n)_{n\geq1}$ \omterm{def:b3:probability:independence}{独立}同分布，其分布函数连续，
称在时刻 $n$ 出现一个\emph{记录}，若
$X_n > \max(X_1, \dots, X_{n-1})$（时刻 $1$ 是
一个记录）。设 $R_n$ 为记录指示变量。
(a) 证明 $\P(R_n = 1) = \frac1n$ \emph{（由对称性，各
$n!$ 种 $X_1, \dots, X_n$ 的排列等可能，
且取值相同的概率为 $0$）}。
(b) 证明 $R_n$ 相互\emph{\omterm{def:b3:probability:independence}{独立}}\emph{（计算与给定记录位置相容的排列数，
或论证
$X_1, \dots, X_{n-1}$ 的相对次序
与 $X_n$ 在其中的秩\omterm{def:b3:probability:independence}{独立}）}。
(c) 由博雷尔–坎泰利引理
（见\cref{thm:b3:probability:borelcantelli}，两部分）可得
几乎必然有无穷多个记录出现，但
相邻时刻 $n, n+1$ 同时出现记录无穷多次的概率
——请判断该概率是 0 还是 1！——并计算
$\sum_n\P(R_n = 1, R_{n+1} = 1)$.
\end{exercise}

\begin{exercise}[$\star\star$]\label{exo:b3:probability:12}
（最长正面游程）掷一枚均匀硬币无穷次，并
记 $L_n$ 为前 $n$ 次投掷中连续正面
构成的最长游程的长度。
(a) 证明：对每个 $\varepsilon > 0$，几乎必然 $L_n \leq
(1 + \varepsilon)\log_2n$ 最终成立 \emph{（
某个长为 $\ell$ 的游程在前 $n$ 次投掷中开始的概率
至多为 $n2^{-\ell}$；沿 $n =
2^k$ 应用博雷尔–坎泰利引理）}。
(b) 证明：几乎必然 $L_n \geq (1 - \varepsilon)\log_2n$
最终成立 \emph{（将前 $n$ 次投掷分为
$\lfloor n/\ell\rfloor$ 个互不相交、长为 $\ell =
\lceil(1 - \varepsilon)\log_2n\rceil$ 的区块，它们
相互\omterm{def:b3:probability:independence}{独立}，每块全为正面的概率为 $2^{-\ell}$，
且无一块全为正面的概率至多为
$\exp(-n2^{-\ell}/\ell)$；沿 $n = 2^k$ 再次求和）}。
(c) 由此得出 $\frac{L_n}{\log_2n} \to 1$ 几乎必然成立：在一百万次
均匀投掷中，应预期出现长约 $20$ 的连续正面游程——
而没有这种游程的数据，很可能是伪造的。
\end{exercise}

\section{问题：埃特马迪的强大数律证明}

\begin{problem}[{周末问题——独立同分布可积变量的
强大数律}]
\label{pb:b3:probability:1}
柯尔莫哥洛夫强大数律——$\frac{S_n}n \to \E X_1$ 几乎必然
对\omterm{def:b3:probability:independence}{独立}同分布的 $X_n \in L^1$ 成立——长期只有繁复的证明；
1981 年，埃特马迪找到一个极为简洁的证明，
只用本章的内容（甚至把\omterm{def:b3:probability:independence}{独立}性减弱为
两两\omterm{def:b3:probability:independence}{独立}）。下面沿用此证明。设 $(X_n)$ 为
两两\omterm{def:b3:probability:independence}{独立}、同分布且\omterm{def:b3:lebesgue:l1}{可积}的序列；
$m = \E X_1$, $S_n = X_1 + \dots + X_n$.

\medskip
\noindent\textbf{第一部分——化归。}
\begin{enumerate}
  \item 证明只需讨论 $X_n \geq 0$ 的情形
        \emph{（将 $X_n = X_n^+ - X_n^-$ 分解，验证两部分
        仍为两两\omterm{def:b3:probability:independence}{独立}、同分布的序列，
        且\omterm{def:b3:lebesgue:l1}{可积}）}。此后假定 $X_n \geq 0$。
  \item （截断）设 $Y_n = X_n\,\mathbf 1_{X_n \leq n}$
        并记 $S_n^* = Y_1 + \dots + Y_n$。证明
        \[
        \sum_{n\geq1}\P(X_n \neq Y_n) =
        \sum_{n\geq1}\P(X_1 > n) \leq \E[X_1] < \infty
        \]
        （见\cref{exo:b3:product:3}），并由
        博雷尔–坎泰利引理推出 $\frac{S_n - S_n^*}{n} \to 0$
        几乎必然。故只需证明 $\frac{S^*_n}n \to m$
        几乎必然。
  \item 证明 $\E Y_n = \E\bigl[X_1\mathbf 1_{X_1\leq
        n}\bigr] \to m$（由单调收敛定理），从而
        $\frac1n\sum_{k\leq n}\E Y_k \to m$（切萨罗平均）：
        只需证明 $\frac{S_n^* - \E S_n^*}{n} \to 0$
        几乎必然。
\end{enumerate}

\medskip
\noindent\textbf{第二部分——方差估计。}
\begin{enumerate}[resume]
  \item 证明
        \[
        \V(Y_n) \leq \E[Y_n^2] = \E\bigl[X_1^2\,\mathbf
        1_{X_1 \leq n}\bigr]
        \]
        并利用层饼公式
        （\cref{prop:b3:product:layercake}），得关键估计
        \[
        \sum_{n\geq1}\frac{\V(Y_n)}{n^2}
        \leq \sum_{n\geq1}\frac1{n^2}\,
        \E\bigl[X_1^2\mathbf 1_{X_1\leq n}\bigr]
        \leq C\,\E[X_1] < \infty
        \]
        \emph{（交换求和与\omterm{def:b3:probability:space}{期望}的顺序——
        级数的托内利定理——并用 $\sum_{n \geq
        x}\frac1{n^2} \leq \frac2{\max(x,1)}$，作内层
        估计 $x^2\sum_{n\geq x}n^{-2} \leq 2x$）}。
\end{enumerate}

\medskip
\noindent\textbf{第三部分——沿几何
子列的收敛性。} 取 $\alpha > 1$，令 $k_j =
\lfloor\alpha^j\rfloor$。
\begin{enumerate}[resume]
  \item 利用两两\omterm{def:b3:probability:independence}{独立}性（方差可加，
        \cref{thm:b3:probability:independence}——验证方差的可加性
        只需两两\omterm{def:b3:probability:independence}{独立}性），
        并由切比雪夫不等式，证明对每个
        $\varepsilon > 0$：
        \[
        \sum_{j\geq1}\P\Bigl(\Bigl|
        \frac{S^*_{k_j} - \E S^*_{k_j}}{k_j}\Bigr| \geq
        \varepsilon\Bigr)
        \leq
        \frac1{\varepsilon^2}\sum_{j\geq1}\frac1{k_j^2}
        \sum_{n\leq k_j}\V(Y_n)
        = \frac1{\varepsilon^2}\sum_{n\geq1}\V(Y_n)
        \sum_{j\,:\,k_j\geq n}\frac1{k_j^2} .
        \]
  \item 证明 $\sum_{j : k_j \geq n}k_j^{-2} \leq
        \frac{C_\alpha}{n^2}$ \emph{（几何级数；
        注意取整：$k_j \geq \frac{\alpha^j}2$ 对
        $\alpha^j \geq 2$ 的情形须小心）}，并由
        第 4 问及博雷尔–坎泰利引理推出：
        \[
        \frac{S^*_{k_j} - \E S^*_{k_j}}{k_j}
        \xrightarrow[j\to\infty]{\text{几乎必然}} 0,
        \qquad\text{因此}\qquad
        \frac{S^*_{k_j}}{k_j} \to m \ \text{几乎必然}
        \]
\end{enumerate}

\medskip
\noindent\textbf{第四部分——夹逼与结论。}
\begin{enumerate}[resume]
  \item 对 $k_j \leq n \leq k_{j+1}$，利用
        $S^*_n$ 的单调性（求和项非负！）证明
        \[
        \frac{k_j}{k_{j+1}}\,\frac{S^*_{k_j}}{k_j}
        \;\leq\; \frac{S^*_n}{n} \;\leq\;
        \frac{k_{j+1}}{k_j}\,\frac{S^*_{k_{j+1}}}{k_{j+1}},
        \]
        并几乎必然推出：
        \[
        \frac m\alpha \leq \liminf\frac{S^*_n}n \leq
        \limsup\frac{S^*_n}n \leq \alpha\,m .
        \]
  \item 沿某一数列令 $\alpha \downarrow 1$，
        得 $\frac{S_n^*}n \to m$ 几乎必然，由（第一部分）
        即\emph{强大数定律}：
        \[
        \boxed{\ \frac{S_n}{n}
        \xrightarrow[n\to\infty]{\text{几乎必然}} \E[X_1].\ }
        \]
  \item 究竟在何处，两两\omterm{def:b3:probability:independence}{独立}性（而非
        完全\omterm{def:b3:probability:independence}{独立}性）已经足够？请列出
        用到\omterm{def:b3:probability:independence}{独立}性类假设的三处。
\end{enumerate}

\medskip
\noindent\textbf{第五部分——附加题。}
\begin{enumerate}[resume]
  \item （博雷尔正规数）证明：
        关于 $\lambda$ 几乎处处的 $x \in \intcc01$ 是
        \emph{在每个基下都是正规数} $b \geq 2$：每个数字
        $0, \dots, b-1$ 出现的渐近频率为
        $\frac1b$ \emph{（固定 $b$ 与某一数字，对指示变量应用
        强大数定律——验证均匀变量的
        $b$ 进制数字是
        \omterm{def:b3:probability:independence}{独立}同分布且在 $\{0,\dots,b-1\}$ 上均匀分布，如
        \cref{thm:b3:probability:existence}——再
        对可数多个概率为一的事件取交
        集）}。给出一个显式的非正规数，
        并思考：该定理断言几乎所有
        数都是正规的，而证明 $\sqrt2$ 或
        $\pi$ 的正规性仍是公开问题。
  \item （有保证的\omterm{ex:b3:probability:sllnapps}{蒙特卡罗}方法）完整论证
        \cref{ex:b3:probability:sllnapps}(b) 中的方法，对
        $g \in L^1(\intcc01^d)$ 的情形：构造
        $\intcc01^d$ 上的\omterm{def:b3:probability:independence}{独立}同分布均匀样本，由
        \cref{thm:b3:probability:existence} 及
        \cref{exo:b3:probability:10}，并说明
        强大数定律给出的结论。
\end{enumerate}

\medskip
\noindent\textbf{第六部分——完全\omterm{def:b3:probability:independence}{独立}性的作用：
极大不等式与随机级数。} 埃特马迪只用到
两两\omterm{def:b3:probability:independence}{独立}性；其余部分则利用
完全（相互）\omterm{def:b3:probability:independence}{独立}的版本。设 $(Z_n)$ 为
$L^2$ 中\omterm{def:b3:probability:independence}{独立}中心化的变量，令 $S_k = Z_1 + \dots + Z_k$（新
记号，与上文的 $X_n$ 无关）。
\begin{enumerate}[resume]
  \item （柯尔莫哥洛夫极大不等式）对 $\varepsilon
        > 0$，证明
        \[
        \P\Bigl(\max_{1\leq k\leq n}\abs{S_k} \geq
        \varepsilon\Bigr) \;\leq\;
        \frac1{\varepsilon^2}\sum_{k=1}^n\V(Z_k) :
        \]
        切比雪夫的代价换来了对极大值的控制\emph{（
        按首个指标 $k$ 划分事件，满足
        $\abs{S_k} \geq \varepsilon$ 的部分；在该部分上写
        $S_n^2 \geq S_k^2 + 2S_k(S_n - S_k)$，并利用
        两组 $(Z_1, \dots, Z_k)$ 与
        $(Z_{k+1}, \dots, Z_n)$ 的\omterm{def:b3:probability:independence}{独立}性，
        \cref{thm:b3:probability:independence}）}。指出
        哪一步上两两\omterm{def:b3:probability:independence}{独立}性将不再
        足够。
  \item （辛钦–柯尔莫哥洛夫单级数定理）推出：
        若 $\sum_n\V(Z_n) < \infty$，则 $\sum_nZ_n$
        几乎必然收敛 \emph{（证明几乎必然地，部分和
        构成柯西列：令 $m \to
        \infty$ 时的极大不等式，作用于
        $Z_{N+1}, \dots, Z_{N+m}$，再令 $N \to
        \infty$）}。
  \item （拉德马赫级数）设 $(\varepsilon_n)$ 为
        \omterm{def:b3:probability:independence}{独立}同分布的符号，$\P(\varepsilon_n = \pm1) =
        \frac12$（\cref{thm:b3:probability:existence}），并令
        $(x_n)$ 为实数列。证明
        $\sum_nx_n\varepsilon_n$ 几乎必然收敛，只要
        $\sum_nx_n^2 < \infty$；并证明，无论
        $(x_n)$ 如何，
        $\sum_nx_n\varepsilon_n$ 收敛的概率为 $0$ 或 $1$
        （\cref{thm:b3:probability:zeroone}）。
  \item 反之，用初等方法。记 $T_n = \sum_{k\leq
        n}x_k\varepsilon_k$ 及 $s_n^2 = \sum_{k\leq
        n}x_k^2$，并假设 $s_n \to \infty$。
        (a) 证明\emph{佩利–齐格蒙德不等式}：对满足
        $Z \geq 0$、$\E Z^2 < \infty$ 与 $0 < \theta <
        1$，
        \[
        \P\bigl(Z > \theta\,\E Z\bigr) \;\geq\; (1 -
        \theta)^2\,\frac{(\E Z)^2}{\E Z^2}
        \]
        \emph{（将 $\E Z$ 在水平 $\theta\E Z$ 处分开，并
        对上段应用柯西–施瓦茨不等式）}。
        (b) 证明 $\E T_n^4 \leq 3s_n^4$。
        (c) 推出 $\P\bigl(\abs{T_n} > \frac{s_n}2\bigr)
        \geq \frac3{16}$，并由此得出
        $\sum_nx_n\varepsilon_n$ 几乎必然发散；因此
        得到二分结论：
        \[
        \sum_nx_n\varepsilon_n\ \text{几乎必然收敛}
        \iff \sum_nx_n^2 < \infty .
        \]
  \item （随机调和级数）推出：
        $\sum_n\frac{\varepsilon_n}{n^s}$ 几乎必然收敛\
        当且仅当 $s > \frac12$。对 $\frac12 < s \leq
        1$ 时，级数几乎必然收敛，而
        $\sum_nn^{-s} = \infty$：随机符号产生
        平方根量级的抵消——可与
        交错级数 $\sum_n\frac{(-1)^n}{n^s}$ 比较，
        它对\emph{每个} $s > 0$ 都收敛。
\end{enumerate}

\medskip
\noindent\textbf{第七部分——集中现象：霍夫丁
不等式。} 强大数定律断言 $\frac{S_n}n \to m$；
而集中不等式说明偏差有多么不可能
\emph{在每个固定的 $n$ 处}。
\begin{enumerate}[resume]
  \item（霍夫丁引理）
        (a) 证明 $\cosh\lambda \leq \eu^{\lambda^2/2}$ 对
        所有 $\lambda \in \R$，通过两个级数的
        逐项比较。
        (b) 设 $Z$ 为中心化\omterm{def:b3:probability:space}{随机变量}，且 $a \leq Z \leq b$，
        $a < b$。证明
        \[
        \E\,\eu^{\lambda Z} \leq
        \exp\Bigl(\frac{\lambda^2(b - a)^2}8\Bigr)
        \]
        \emph{（将 $\eu^{\lambda z}$ 在 $\intcc ab$ 上用其弦控制，
        取\omterm{def:b3:probability:space}{期望}，并研究 $\varphi(t)
        = -pt + \log(1 - p + p\eu^t)$，其中 $p =
        \frac{-a}{b-a}$、$t = \lambda(b - a)$，证明
        $\varphi(0) = \varphi'(0) = 0$ 且 $\varphi''
        \leq \frac14$）}。
  \item（霍夫丁不等式）设 $X_1, \dots, X_n$ 为
        \omterm{def:b3:probability:independence}{独立}，满足 $a_i \leq X_i \leq b_i$，且 $S_n =
        X_1 + \dots + X_n$。证明：对 $t > 0$，
        \[
        \P\bigl(S_n - \E S_n \geq t\bigr) \leq
        \exp\Bigl(\frac{-2t^2}{\sum_{i=1}^n(b_i -
        a_i)^2}\Bigr),
        \]
        对下尾也有相同的界，
        \emph{（指数型切比雪夫：界定
        $\E\,\eu^{\lambda(S_n - \E S_n)}$，利用
        \omterm{def:b3:probability:independence}{独立}性及第 17 题，然后对
        $\lambda > 0$ 优化）}。
  \item（强大数定律，有界情形，附收敛速度）设
        $X_i$ 为取值于 $\intcc ab$ 中的\omterm{def:b3:probability:independence}{独立}同分布\omterm{def:b3:probability:space}{随机变量}，且 $m
        = \E X_1$。证明
        \[
        \P\Bigl(\Bigl|\frac{S_n}n - m\Bigr| \geq
        \varepsilon\Bigr) \leq
        2\exp\Bigl(\frac{-2n\varepsilon^2}{(b - a)^2}\Bigr)
        \]
        并推出 $\frac{S_n}n \to m$ 几乎必然成立，这借助
        博雷尔–坎泰利引理，得到有界变量强大数定律的第二种证明 ——
        不作截断，且
        在每个有限 $n$ 处有指数速度，但要求加项有界
        且完全\omterm{def:b3:probability:independence}{独立}。比较这些假设
        与埃特马迪的假设。
  \item（\omterm{ex:b3:probability:sllnapps}{蒙特卡罗}，对固定的 $n$ 有保证）设 $g
        \colon \intcc01^d \to \intcc01$ 为\omterm{def:b3:lebesgue:measurable}{可测函数}，且
        $(U_k)$ 为第 11 题中的\omterm{def:b3:probability:independence}{独立}同分布均匀样本。
        给定 $\varepsilon, \delta > 0$，证明
        \[
        n \geq \frac{\log(2/\delta)}{2\varepsilon^2}
        \implies
        \P\Bigl(\Bigl|\frac1n\sum_{k=1}^ng(U_k) -
        \int g\,\dd\lambda_d\Bigr| \geq \varepsilon\Bigr)
        \leq \delta,
        \]
        并计算当 $\varepsilon =
        \delta = 10^{-2}$ 时的门槛值。该界不依赖 $d$：
        请与第 11 题以及确定性
        网格作比较。
\end{enumerate}

\medskip
\noindent\textbf{第八部分——随机游走有多大？
迈向重对数律。} 设 $S_n = \varepsilon_1 +
\dots + \varepsilon_n$ 为由
        \omterm{def:b3:probability:independence}{独立}同分布的公平符号构成的简单随机游走。
\begin{enumerate}[resume]
  \item（次高斯尾）证明 $\E\,\eu^{\lambda S_n} =
        (\cosh\lambda)^n \leq \eu^{n\lambda^2/2}$，并
        对 $x > 0$，由此推出：
        \[
        \P(S_n \geq x) \leq \eu^{-x^2/(2n)},
        \qquad
        \P(\abs{S_n} \geq x) \leq 2\,\eu^{-x^2/(2n)} .
        \]
  \item 借助博雷尔–坎泰利引理，推出
        \[
        \limsup_{n\to\infty}\frac{\abs{S_n}}
        {\sqrt{2n\log n}} \leq 1 \quad\text{几乎必然}
        \]
        \emph{（对 $\eta > 0$，在 $x =
        (1 + \eta)\sqrt{2n\log n}$ 处对尾概率界求和，再对
        $\eta = \frac1p$ 取交）}。特别地，随机游走
        位于中心极限定理尺度 $\sqrt n$ 上，只差一个对数因子
        ——远小于粗糙的界 $\abs{S_n} \leq n$。
  \item 沿倍增子列 $n_j = 2^j$，证明
        \[
        \limsup_{j\to\infty}\frac{S_{n_j}}
        {\sqrt{2n_j\log\log n_j}} \leq 1 \quad\text{几乎必然},
        \]
        并加以说明：\emph{重对数律
        }（辛钦；哈特曼–温特纳，对一般中心化的 $L^2$ 加项）
        断言：
        \[
        \limsup_{n\to\infty}\frac{S_n}
        {\sqrt{2n\log\log n}} = 1 \quad\text{几乎必然}
        \]
        请准确说明，刚证明的子列估计与本断言的上半部分
        究竟有何区别。
        （须控制每个块内的 $\max_{n_j \leq n \leq
        n_{j+1}}S_n$，这需要
        在\emph{指数}尺度上的极大不等式）。
        并定量检验第 12 题的不等式对此目的
        是否太弱。下半部分
        依赖第二博雷尔–坎泰利引理，
        并作用于\omterm{def:b3:probability:independence}{独立}的块上；两部分都是
        本科三年级的正规内容，
        宜在专门的概率论课程中讲授。

  \item （有限类上的一致偏差）设 $A_1,
        \dots, A_N$ 为可重复试验中的事件，
        用经验频率估计每个概率，
        其中 $\hat p_i$ 由 $n$ 次\omterm{def:b3:probability:independence}{独立}同分布重复试验得到。
        结合霍夫丁不等式与并集界，
        证明
        \[
        \P\Bigl(\max_{i\leq N}\,\abs{\hat p_i - \P(A_i)} >
        \varepsilon\Bigr) \;\leq\; 2N\,\eu^{-2n\varepsilon^2},
        \]
        并推出样本量规则：$n \geq
        \frac{\ln(2N/\delta)}{2\varepsilon^2}$ 可保证
        所有 $N$ 个估计同时满足
        $\varepsilon$ 精度，概率至少为 $\geq 1 -
        \delta$。计算 $n$，取 $N = 10^6$、$\varepsilon =
        0.01$、$\delta = 0.05$ 时的样本量，此即一致性的
        对数代价。
  \item （随机调和窗口）结合
        随机级数理论的两半，证明：
        对\omterm{def:b3:probability:independence}{独立}同分布的符号 $(\varepsilon_n)$，级数
        $\sum_n\frac{\varepsilon_n}{n^\alpha}$ 在
        $\alpha > \frac12$ 时几乎必然收敛，在
        $\alpha \leq \frac12$ 时几乎必然发散。与绝对收敛
        （要求 $\alpha > 1$）相比：在
        窗口 $\alpha \in \intoc{\frac12}1$ 上，收敛
        是真正的概率现象——
        源于抵消，而非项的大小。
\end{enumerate}
\end{problem}
